% !TeX root = ../sections/02_exercises.tex
\subsection{2-B-8.}
\begin{exercise}
	数列 \(\{a_n\}_{n\in\mathbb{N}}\) が，ある \(0<r<1\) に対して
	\[
	|a_{n+2}-a_{n+1}|
	\leq
	r|a_{n+1}-a_n|
	\qquad
	(n=1,2,3,\ldots)
	\]
	を満たすとき，数列 \(\{a_n\}_{n\in\mathbb{N}}\) は収束することを証明せよ。
\end{exercise}

\begin{background}
	この問題では，隣り合う項の差が一定の割合 \(r\) で小さくなることを利用する。
	ここで
	\[
	0<r<1
	\]
	であるから，差の大きさは等比数列のように減少する。
	
	実際，仮定から
	\[
	|a_{n+1}-a_n|
	\leq
	r^{n-1}|a_2-a_1|
	\]
	が成り立つことが分かる。
	
	そして，\(m>n\) のとき，
	\[
	a_m-a_n
	=
	(a_m-a_{m-1})
	+\cdots+
	(a_{n+1}-a_n)
	\]
	と分解できる。
	
	三角不等式と等比級数の評価を使えば，
	\[
	|a_m-a_n|
	\]
	を十分小さくできる。
	したがって \(\{a_n\}\) は Cauchy 列であり，
	\(\mathbb{R}\) の完備性より収束する。
\end{background}

\begin{strategy}
	まず，帰納法により
	\[
	|a_{n+1}-a_n|
	\leq
	r^{n-1}|a_2-a_1|
	\]
	を示す。
	
	\(n=1\) のときは
	\[
	|a_2-a_1|
	\leq
	r^0|a_2-a_1|
	\]
	なので成り立つ。
	
	いま
	\[
	|a_{n+1}-a_n|
	\leq
	r^{n-1}|a_2-a_1|
	\]
	が成り立つとする。
	仮定より
	\[
	\begin{aligned}
		|a_{n+2}-a_{n+1}|
		&\leq
		r|a_{n+1}-a_n| \\
		&\leq
		r\cdot r^{n-1}|a_2-a_1| \\
		&=
		r^n|a_2-a_1|.
	\end{aligned}
	\]
	よって帰納法により，
	\[
	|a_{n+1}-a_n|
	\leq
	r^{n-1}|a_2-a_1|
	\]
	がすべての \(n\) で成り立つ。
	
	次に，\(\{a_n\}\) が Cauchy 列であることを示す。
	\(m>n\) とする。
	このとき
	\[
	a_m-a_n
	=
	\sum_{k=n}^{m-1}(a_{k+1}-a_k)
	\]
	であるから，三角不等式より
	\[
	\begin{aligned}
		|a_m-a_n|
		&\leq
		\sum_{k=n}^{m-1}|a_{k+1}-a_k| \\
		&\leq
		\sum_{k=n}^{m-1}
		r^{k-1}|a_2-a_1| \\
		&\leq
		|a_2-a_1|
		\sum_{k=n}^{\infty}r^{k-1}.
	\end{aligned}
	\]
	
	等比級数の和より，
	\[
	\sum_{k=n}^{\infty}r^{k-1}
	=
	\frac{r^{n-1}}{1-r}.
	\]
	したがって
	\[
	|a_m-a_n|
	\leq
	\frac{|a_2-a_1|}{1-r}r^{n-1}.
	\]
	
	右辺は \(n\to\infty\) で \(0\) に収束する。
	よって，任意の \(\varepsilon>0\) に対して，
	ある \(N\in\mathbb{N}\) が存在して，
	\(m>n\geq N\) ならば
	\[
	|a_m-a_n|<\varepsilon
	\]
	となる。
	
	したがって，\(\{a_n\}\) は Cauchy 列である。
	実数列の Cauchy 列は収束するので，
	\[
	\{a_n\}
	\]
	は収束する。
\end{strategy}